\documentclass[11pt,a4paper]{article}
\usepackage[T1]{fontenc}
\usepackage[utf8]{inputenc}
\usepackage[english]{babel}
\usepackage{lmodern,microtype,amsmath,amssymb,amsthm,mathtools}
\usepackage[margin=25mm]{geometry}
\usepackage{booktabs,enumitem,xcolor}
\usepackage[breakable]{tcolorbox}
\newtcolorbox{notes}{breakable,
  colback=yellow!8!white,colframe=orange!60!black,
  colbacktitle=yellow!25!white,coltitle=black,
  title=Notes,fonttitle=\bfseries,boxrule=0.7pt,arc=2mm}
\usepackage[colorlinks=true,linkcolor=teal!65!black,urlcolor=teal!65!black,citecolor=teal!65!black]{hyperref}
\newcommand{\dep}[2]{\mathord{=}(#1;#2)}
\newcommand{\con}[1]{\mathord{=}(#1)}
\newcommand{\ND}[3]{#1\xrightarrow{\leq #3}#2}
\newcommand{\fin}{\mathrm{fin}}
\newcommand{\RE}{\mathsf{r.e.}}
\newcommand{\coRE}{\mathsf{co\text{-}r.e.}}
\newtheorem{theorem}{Theorem}[section]
\newtheorem{lemma}[theorem]{Lemma}
\newtheorem{corollary}[theorem]{Corollary}
\theoremstyle{definition}
\newtheorem{definition}[theorem]{Definition}
\setlist{itemsep=3pt,topsep=5pt}
\title{Undecidability of Implication for\\Disjunctions of Dependence Atoms}
\author{}
\date{24 October 2026}
\begin{document}
\maketitle
\vspace{-8mm}
\begin{notes}
\textbf{Miika Hannula (University of Tartu):} This document is fully AI written. After iterative improvements on the lower bound, I suggested to investigate some kind of small model property for decidability, and if this fails to try identify some kind of obstacle. This led to the proof, after which I   requested a couple minor modifications.
\tcblower
\small
The research note was generated using ChatGPT (OpenAI, GPT-6 Astra,
Ultra mode) under the ChatGPT for Academic Researchers complimentary account.
\end{notes}

\begin{abstract}
We consider implication from a finite set of disjunctions of dependence
atoms to one such disjunction, under team semantics and
without a bound on the value domain. We prove that finite and unrestricted
implication coincide and that their common decision problem is
$\coRE$-complete. Hardness already holds for a fixed consequent with six
atom occurrences, with at most five occurrences in each premise, and with
only five variables exempt from explicit three-value bounds. The proof
reduces two-counter machine halting to the existence of a finite
counterteam. A Hall obstruction forces the complete incidence structure
needed for the simulation; two matching channels then enforce control
flow and exact counter updates. A three-marker construction compiles
Boolean row conditions and conditional dependence requirements into the required
pure dependence language. Consequently, every failure of implication has
a finite counterteam, but no computable function of the input bounds the
size of a smallest one.
\end{abstract}

\section{The implication problem}
We follow the standard team semantics of
V\"a\"an\"anen~\cite[Chapter~3]{Vaananen07}.

\begin{definition}[Teams and dependence atoms]
\label{def:dependence}
Let $V$ be a finite set of variables and $D$ a nonempty value domain.
An \emph{assignment} is a function $s\colon V\to D$, and a
\emph{team} is a possibly empty set $T\subseteq D^V$ of assignments.
For $X\subseteq V$ and $y\in V$, the \emph{dependence atom}
$\dep Xy$ states that the values of the variables in $X$ determine
the value of $y$. Precisely,
\[
 T\models\dep Xy
 \quad\Longleftrightarrow\quad
 \forall s,t\in T\;
 \bigl[(\forall x\in X\;s(x)=t(x))\Longrightarrow s(y)=t(y)\bigr].
\]
For $X=\{x_1,\ldots,x_m\}$ we also write $\dep{x_1\cdots x_m}{y}$;
the order of the determinant variables is immaterial.
We write $\con y$ for the constancy atom $\dep\varnothing y$.
It holds exactly when $y$ has at most one value in $T$.
\end{definition}

A \emph{dependence disjunction} is a nonempty finite formula
\[
 \psi=\bigvee_{i=1}^k\dep{X_i}{y_i},\qquad k\geq1.
\]
Its semantics is
\[
 T\models\psi
 \quad\Longleftrightarrow\quad
 \exists T_1,\ldots,T_k\subseteq T:
 \quad T=\bigcup_iT_i\ \text{ and }\ T_i\models\dep{X_i}{y_i}.
\]
The parts may overlap and may be empty.\footnote{Requiring pairwise
disjoint parts gives the same semantics: replace $T_i$ by
$T_i\setminus\bigcup_{j<i}T_j$. Each smaller team still satisfies its
dependence atom. Thus strict and lax disjunction agree in this fragment.}
Repeated atom occurrences are
retained: for example, $\con y\vee\con y$ allows two distinct $y$-values.
There are no other atoms, constants, relation symbols, or first-order
literals in the input language.

For a set $\Sigma$ of formulas, $T\models\Sigma$ means that
$T\models\psi$ for every $\psi\in\Sigma$.

\begin{center}
\setlength{\fboxsep}{9pt}
\fbox{\begin{minipage}{\dimexpr\linewidth-2\fboxsep-2\fboxrule\relax}
\textbf{Implication for dependence disjunctions}
\par\medskip
\textbf{Input.}\enspace A finite set $\Sigma$ of dependence disjunctions
and one dependence disjunction $\varphi$. Let $V$ be the set of variables
occurring in $\Sigma\cup\{\varphi\}$. All variables and atom occurrences
are listed explicitly.
\par\medskip
\textbf{Question.}\enspace Does $\Sigma\models\varphi$ hold? That is,
is the following condition true?
\[
 \forall D\ne\varnothing\;\forall T\subseteq D^V:\qquad
 T\models\Sigma\ \Longrightarrow\ T\models\varphi.
\]
\textbf{Finite-team variant.}\enspace Ask the same question with $T$
required to be finite; write $\Sigma\models_{\fin}\varphi$.
In both variants the value domain $D$ is unrestricted.
\end{minipage}}
\end{center}

All equality tests compare values in the \emph{same} variable; independently
renaming the values of different variables does not change satisfaction.

\begin{samepage}
\begin{theorem}[Main result]
\label{thm:main}
Finite and unrestricted implication coincide. Their common decision
problem is $\coRE$-complete under computable many-one reductions.
Hardness holds already for the fixed consequent
\[
 \varphi_*=\dep wz\vee\dep wz\vee
             \dep Uw\vee\dep Uw\vee
             \dep Vw\vee\dep Vw.
\]
Each premise has at most five atom occurrences. Moreover, every variable
other than $U,V,X,Y,w$ can be explicitly restricted to at most three
values.
\end{theorem}
\end{samepage}

Here $\RE$ denotes the recursively enumerable sets and $\coRE$ their
complements. The hardness proof uses two-counter machine halting
\cite{Dudenhefner22}. Sections~\ref{sec:compiler}--\ref{sec:reduction}
give the reduction, including the compilation into pure dependence
formulas. We first establish finite counterteams and the upper bound.

Implication for this disjunction fragment was considered in
Kontinen's thesis \cite[Section~7.3]{Kontinen10}. The more restricted
case involving disjunctions of two dependence atoms is raised as a
further question by Fr\"ohlich, Kolaitis, and Meier
\cite[Section~6]{FKM26}. The present reduction uses longer disjunctions
and does not settle that two-occurrence restriction. Our argument is
self-contained apart from propositional compactness, Hall's theorem,
and the undecidability of two-counter machine halting.

\section{Finite counterteams and enumeration}
All formulas in the language are downward closed: if $S\subseteq T$ and
$T\models\psi$, then $S\models\psi$. For an atom this is immediate;
for a disjunction, intersect each part of a cover with $S$.

\begin{lemma}[Finite character]
\label{lem:finite-character}
For every formula $\psi$ and team $T$,
\[
 T\models\psi
 \quad\Longleftrightarrow\quad
 \text{every finite }S\subseteq T\text{ satisfies }\psi.
\]
\end{lemma}
\begin{proof}
Only the reverse implication needs proof. Write
$\psi=\bigvee_{i=1}^k\dep{X_i}{y_i}$. Introduce propositional variables
$p_{s,i}$ for $s\in T$ and impose
\[
 \bigvee_{i=1}^k p_{s,i}\quad(s\in T),
 \qquad
 \neg p_{s,i}\vee\neg p_{t,i}
 \quad\bigl(s|_{X_i}=t|_{X_i},\ s(y_i)\ne t(y_i)\bigr).
\]
Every finite subset of these clauses mentions finitely many rows, and
a cover of that finite subteam satisfies the clauses. Propositional
compactness supplies a model of all clauses. Its sets
$T_i=\{s:p_{s,i}\text{ is true}\}$ form the required cover.
\end{proof}

\begin{corollary}
\label{cor:finite}
$\Sigma\models_{\fin}\varphi$ if and only if $\Sigma\models\varphi$.
\end{corollary}
\begin{proof}
If $T\models\Sigma$ but $T\not\models\varphi$, the lemma gives a
finite $S\subseteq T$ with $S\not\models\varphi$. Downward closure
preserves all premises on $S$.
\end{proof}

\begin{lemma}
\label{lem:enumeration}
Nonimplication is recursively enumerable.
\end{lemma}
\begin{proof}
Enumerate $r=1,2,\ldots$ and all teams on the input variables with
values in $[r]$ and at most $r$ rows. Every finite team appears up to
independent renaming within each variable. For a given finite team,
satisfaction of each formula is decidable by enumerating its row
colourings. Stop upon finding a model of $\Sigma$ that violates
$\varphi$. Corollary~\ref{cor:finite} ensures termination on every
nonimplication instance.
\end{proof}

We will also use projection invariance. If a formula uses only variables
in $A$, it holds in $T$ exactly when it holds in $\pi_A(T)$: covers
project to the latter and lift back to the former. Injective renaming
within each variable preserves all atoms and their disjunctions.

\section{Compiling Boolean conditions into dependence disjunctions}
\label{sec:compiler}
Write $\ND Xyr$ for a disjunction of $r$ copies of $\dep Xy$.
It states that each $X$-fibre has at most $r$ distinct $y$-values:
each part permits one value, and conversely those values can be assigned
to $r$ parts separately in every fibre. This is an abbreviation,
not an additional atom type.

\subsection{Filtering complete three-marker fibres}
Let $A$ be a set of \emph{semantic variables}, and let $z\notin A$
be a new marker variable. Under the premise $\ND\varnothing z3$, define
\[
 C_A(T)=\{a\in\pi_A(T):
       |\{s(z):s\in T,\ s|_A=a\}|=3\}.
\]
This is the \emph{complete-fibre core}, a team on $A$. Auxiliary
variables of $T$ are ignored by this projection.

\begin{lemma}[Filtering]
\label{lem:filter}
If $D$ uses only variables in $A$, then
\[
 T\models\dep Az\vee\dep Az\vee D
 \quad\Longleftrightarrow\quad C_A(T)\models D.
\]
\end{lemma}
\begin{proof}
Each $\dep Az$-part covers at most one marker value per fibre.
Thus at least one row in every complete fibre belongs to the parts
for $D$. Their projections cover $C_A(T)$.
Conversely, the first two parts can cover every incomplete fibre and
two marker layers of every complete fibre. Lift a $D$-cover of the
core to all remaining rows.
\end{proof}

\subsection{Equality pairs encode three colours}
For an ordinary graph-vertex variable $c$, impose
\begin{equation}
 \ND\varnothing c2,\qquad \dep{Az}c,\qquad \ND{Ac}z2.
 \label{eq:node}
\end{equation}
On a complete fibre, the three marker rows have exactly two $c$-values.
They therefore select one of the three possible equal pairs. This pair
is the vertex's colour. An edge between vertex variables $c,d$ is
encoded by
\begin{equation}
 \dep{Acd}z.
 \label{eq:edge}
\end{equation}
On a complete fibre this holds exactly when their equal pairs differ:
an identical pair violates the atom, whereas two different pairs have
only singleton common equivalence classes.

Take anchor variables $c_F,c_T,c_B$, give each the conditions
\eqref{eq:node}, impose their triangle using \eqref{eq:edge}, and add
$\dep z{c_F}$, $\dep z{c_T}$, and $\dep z{c_B}$.
If the core is nonempty, their three different pairs are independent
of the semantic fibre. Rename the three markers so that
\[
 F=\{1,2\},\qquad T=\{1,3\},\qquad B=\{2,3\}.
\]
An edge to the anchor $B$ restricts a vertex to the Boolean colours
$F,T$. This encoding never compares values of different variables.

\subsection{A consistent interpretation of semantic Boolean variables}
Let $\mathcal B\subseteq A$ be the variables intended to be Boolean.
For each $v\in\mathcal B$, impose $\ND\varnothing v2$ and introduce
a switch variable $w_v$ with
\begin{align}
 &\ND\varnothing{w_v}3,\quad \dep{vz}{w_v},\quad
   \ND v{w_v}2,\quad \ND{v w_v}z2,\label{eq:switch-a}\\
 &\dep{c_Fw_v}v,\qquad \dep{c_Tw_v}v,\qquad
   \dep{A c_Bw_v}z.\label{eq:switch-b}
\end{align}
Switches do not receive the ordinary vertex's global two-value bound.
The variables in $A\setminus\mathcal B$ remain unrestricted.

\begin{lemma}[Boolean decoding of a partial team]
\label{lem:decode}
If the core is nonempty, each $v\in\mathcal B$ has a globally
consistent injective interpretation
\[
 \eta_v:\operatorname{rng}_{C_A(T)}(v)\longrightarrow\{0,1\},
\]
given by the switch colour $F$ or $T$. Conversely, all switch
conditions admit a canonical realization for any prescribed semantic
Boolean values.
\end{lemma}
\begin{proof}
On a complete fibre, $w_v$ has exactly two values: the bound
$\ND v{w_v}2$ gives at most two, while $\ND{v w_v}z2$ excludes a
constant value on all three markers. The edge to $B$ makes its colour
Boolean. The atom $\dep{vz}{w_v}$ makes this marker triple depend
only on the actual value of $v$.

Suppose two distinct $v$-values occur in the core. Their two-element
sets of switch values intersect, because the switch has at most three
values globally. If a shared value occurs at markers $r,s$ in the
two fibres, the atoms $\dep{c_Fw_v}v$ and $\dep{c_Tw_v}v$ require
both $c_F(r)\ne c_F(s)$ and $c_T(r)\ne c_T(s)$. For the fixed
anchor pairs this is possible only for $(r,s)=(2,3)$ or $(3,2)$.
The shared value must therefore occur as a singleton in each triple,
at opposite markers. The remaining equal pairs are the opposite
Boolean colours $F,T$. This proves injectivity. Values occurring
only in incomplete fibres need no interpretation.

For the canonical realization, use the anchor triples
\[
 c_F=(0,0,1),\quad c_T=(0,1,0),\quad c_B=(0,1,1),
\]
and the switch triples
\[
 w_v=\begin{cases}(0,0,1),&v=0,\\(2,1,2),&v=1.\end{cases}
\]
Direct inspection verifies all conditions. These numbers are values
in a constructed team, not constant symbols in the formulas.
\end{proof}

\subsection{Arbitrary Boolean row conditions}
A finite Boolean circuit can be converted into a polynomial-size
3CNF formula with auxiliary variables, preserving the extendability
of each input assignment. For example, a gate $h=a\wedge b$ is
expressed by $(\neg h\vee a)$, $(\neg h\vee b)$, and
$(h\vee\neg a\vee\neg b)$. Negation and disjunction have analogous
constant-size encodings; the output is required to be true. Repeating
literals pads shorter clauses to length three.

Here is an explicit colouring graph for the resulting 3CNF. Each
Boolean variable has vertices $v,\bar v$ with edges
$v\bar v,vB,\bar vB$, forcing opposite Boolean colours. For a semantic
bit, use its switch as the vertex $v$; all other vertices are ordinary
auxiliary variables.

The condition $\operatorname{NAE}(p,q,r)$, meaning that three Boolean
inputs are not all equal, is encoded by a new triangle $t_1,t_2,t_3$
and the edges $t_1p,t_2q,t_3r$. Equal inputs leave only two colours
for the triangle. For unequal inputs, assign the opposite Boolean
colour and $B$ to the two vertices whose inputs agree, and the other
Boolean colour to the third vertex. This is a proper colouring.

For a clause $\ell_1\vee\ell_2\vee\ell_3$, introduce a new Boolean
vertex $s$ with its complement and impose two such gadgets:
\[
 \operatorname{NAE}(\ell_1,\ell_2,s),\qquad
 \operatorname{NAE}(\bar s,\ell_3,F).
\]
If all literals are false, the first condition forces $s=T$ and the
second fails. If $\ell_3=T$, the second is automatic and $s$ can
satisfy the first. If $\ell_3=F$ but one of the other literals is true,
choose $s=F$. Thus the graph extends a prescribed input colouring
exactly when the clause holds.

Encode vertices and edges using \eqref{eq:node}--\eqref{eq:edge}.
On each complete fibre, the resulting pure dependence premises force
the desired Boolean row condition. Conversely, each satisfying input
assignment extends to a local graph colouring. Give every ordinary
vertex the canonical triple $c_F,c_T,c_B$ for its colour. All auxiliary
variables then have at most three values globally, regardless of the
number of semantic fibres.

\subsection{Conditional dependence and the compilation lemma}
\begin{lemma}[Conditional dependence]
\label{lem:conditional}
Let $b\in\mathcal B$ and let $D$ use only semantic variables in $A$.
The premise
\[
 \dep{A c_T w_b}z\vee D
\]
forces $D$ on the subteam of the core where $b$ is decoded as true.
Conversely, the canonical three-marker expansion satisfies this
premise whenever that selected semantic subteam satisfies $D$.
\end{lemma}
\begin{proof}
On a complete fibre with $b$ true, $c_T$ and $w_b$ have the same
pair. The first part cannot cover both rows of that pair, so at least
one row belongs to the parts for $D$. Projection gives a $D$-cover
of the selected core.
Conversely, lift a $D$-cover to all true-$b$ fibres. On a false-$b$
fibre, $c_T$ and $w_b$ have different pairs, so $(c_T,w_b)$ distinguishes
all three markers. Assign all these fibres to the first part. The
key $A$ prevents interactions between different fibres.
\end{proof}

\begin{lemma}[Compilation]
\label{lem:compiler}
Consider a finite semantic theory with variables $A$, Boolean variables
$\mathcal B\subseteq A$, a Boolean row condition $\Gamma$, ordinary
dependence disjunctions, and conditions of the form
``the subteam selected by the Boolean condition $P$ satisfies the
dependence atom $\alpha$''. One can effectively construct a pure dependence
 disjunction theory $\widehat\Sigma$ such that:
\begin{enumerate}
\item the nonempty core of every $T\models\widehat\Sigma$ satisfies
      the semantic theory after Boolean decoding;
\item every finite semantic model has a finite canonical expansion
      satisfying $\widehat\Sigma$, with its prescribed Boolean guard
      extension as the core;
\item all new variables have explicit bounds of three values.
      In addition to the original disjunctions, at most three atom
      occurrences per premise are needed; a conditional requirement uses two.
\end{enumerate}
\end{lemma}
\begin{proof}
For each condition $P$, add a semantic Boolean variable $b_P$ and
include $b_P\leftrightarrow P$ in the row condition. From this point,
$A$ includes these guard variables. Use the colouring graph above for
the combined row condition, switches for every semantic Boolean bit,
and the marker bound $\ND\varnothing z3$. Retain the ordinary disjunctions
and compile each conditional requirement by Lemma~\ref{lem:conditional}.

The first assertion follows from projection invariance, injectivity
of Boolean decoding, and the constructions above. For the second,
give each semantic row three markers, canonical anchors and switches,
and a graph colouring extending its satisfying Boolean assignment.
Lift ordinary covers to all three markers; use the preceding lemma
for conditional covers. All auxiliary values are in $\{0,1\}$ or
$\{0,1,2\}$, so explicit three-value bounds can be added.
The construction has polynomial size in the input circuit description.
\end{proof}

\section{A Hall obstruction forces complete ports}
Write $k\cdot\alpha$ for a disjunction of $k$ copies of $\alpha$.
First work on a finite semantic team $R$ whose key is $w$: no two
rows have the same $w$-value. View its rows as the edges of a bipartite
multigraph. The variables $U,V$ specify the endpoints, and $w$
distinguishes parallel edges. Put
\begin{align}
 P_U&=(4\cdot\dep Uw)\vee\con w,\label{eq:degree-u}\\
 P_V&=(4\cdot\dep Vw)\vee\con w,\label{eq:degree-v}\\
 D&=(2\cdot\dep Uw)\vee(2\cdot\dep Vw).\label{eq:core-target}
\end{align}

\begin{lemma}[Capacity-two Hall obstruction]
\label{lem:hall}
If $R\models P_U,P_V$ but $R\not\models D$, some subteam
$S\subseteq R$ has exactly one $U$-vertex and one $V$-vertex of
degree five, with every other occurring vertex of degree four.
Moreover,
\[
 |S|=2\bigl(|\pi_U(S)|+|\pi_V(S)|\bigr)+1.
\]
\end{lemma}
\begin{proof}
A $D$-cover is equivalent to assigning each edge to one of two
slots at its $U$-endpoint or one of two slots at its $V$-endpoint,
without reusing a slot. Indeed, each of the four parts allows at most
one edge at an endpoint of the corresponding sort. Hall's theorem
\cite{Hall35} gives an edge set $S$ with
$|S|>2(|\pi_U(S)|+|\pi_V(S)|)$.

Since $w$ is a key, the constancy part of $P_U$ contains at most one
edge. Hence at most one $U$-vertex has degree five, and all others have
degree at most four. The same holds on the $V$-side and is preserved
in $S$. Writing $n=|\pi_U(S)|+|\pi_V(S)|$, we obtain
\[
 4n+2\ \leq\ 2|S|\ =\sum_v\deg_S(v)\ \leq\ 4n+2.
\]
Every inequality is equality, proving the assertion.
\end{proof}

\section{The semantic counter-machine construction}
\label{sec:machine}
Halting for two-counter machines is $\RE$-complete: Turing machine
halting many-one reduces to it \cite[Lemma~3]{Dudenhefner22}.
We use the following normal form. Both counters start at zero and are
cleared before halting. Any prescribed initial counter values can be
built by a finite initialization program, and final values can be
cleared by decrement/zero-test loops. The start state is fresh and has
no incoming ordinary transition; the halt state is fresh and has no
outgoing one. Duplicate states to record time parity. The start state
has $U$-parity and the halt state $V$-parity; insert final no-op steps
as necessary. Fresh intermediate states make each allowed ordered
pair of states determine a unique operation and zero-test guard.
The operations are increment, positive decrement, zero test, and
no-op. A conditional zero/decrement instruction is represented by
transitions to two distinct intermediate states. All transformations
are effective and preserve halting.

The unbounded semantic variables are $U,V,X,Y,w$. The remaining
semantic data
\[
 t,k_U,k_V,q_U,q_V,\delta
\]
are finite-valued codes represented by Boolean bits. Here $X,Y$
name configurations, $U,V$ name their tokens, and $q_U,q_V$ are
control states. A token kind is one of
\[
 S,\ C_1,\ C_2,\ B,
\]
meaning control token, a unit of the first or second counter, or a
blank. Edge types are $a,b,c,d,e_1,e_2,e_3$. Their direction code is
$\delta=0$ for $a,c$, $\delta=1$ for $b,d$, and $\delta=2$ for $e_i$.
Require valid codes throughout, with $q_U,q_V$ belonging to the
appropriate parity classes of states.

All conditions below are imposed on the semantic team. Determinant
codes denote all their bits; output codes abbreviate one dependence
atom per bit.
Include the two degree premises
\eqref{eq:degree-u}--\eqref{eq:degree-v} and the following conditions.
\begin{enumerate}[label=\textup{(\arabic*)}]
\item The variable $w$ determines every semantic variable, including
      all Boolean guard variables introduced by compilation. Impose
\[
 \begin{gathered}
 \dep U X,\quad \dep U{k_U},\quad \dep U{q_U},\quad
 \dep V Y,\quad \dep V{k_V},\quad \dep V{q_V},\\
 \dep X{q_U},\qquad \dep Y{q_V},\qquad
 \dep{Ut}w,\qquad\dep{Vt}w.
 \end{gathered}
\]
Thus token data are consistent, and every token has at most one
incident edge of each type.
\item On rows with $k_U=S$, impose $\dep X U$; on rows with $k_V=S$,
impose $\dep Y V$. Thus every configuration has at most one control
token. On rows with $q_U=\mathrm{start}$ impose $\con X$, and on
rows with $q_V=\mathrm{halt}$ impose $\con Y$. The start and halt
configurations, if present, are globally unique. Only kinds $S,B$
are allowed in these configurations.
\item An $e_i$-edge is allowed only between the start configuration's
$U$-control token and the halt configuration's $V$-control token.
These two tokens forbid types $b,d$. Every other token permits only
types $a,b,c,d$.
\item Impose the dependence atoms
\[
 \dep{X\delta}Y,\qquad \dep{Y\delta}X.
\]
For types $a,c$, the directed transition goes from $(X,q_U)$ to
$(Y,q_V)$; for $b,d$, from $(Y,q_V)$ to $(X,q_U)$. The ordered state
pair must be an allowed machine transition. The sole exception is
a \emph{reset} from halt to start on $b,d$-edges with both endpoint
kinds $B$. No transition condition applies to the $e_i$-edges.
\item The first channel, types $a,b$, requires
$k_U=S\leftrightarrow k_V=S$. It does not otherwise restrict
pairings between noncontrol token kinds.
\item On the second channel, types $c,d$, the pair $(L,R)$ of source
and destination kinds follows the table below. Reset edges allow
only $(B,B)$. A zero test of counter $j$ additionally forbids source
kind $C_j$.
\end{enumerate}
\begin{center}
\small
\begin{tabular}{@{}ll@{}}
\toprule
Operation & Allowed pairs $(L,R)$\\
\midrule
No-op or zero test & $(S,S),(B,B),(C_1,C_1),(C_2,C_2)$\\
Increment $j$ & $(S,C_j),(C_j,C_j),(C_{3-j},C_{3-j}),(B,B),(B,S)$\\
Decrement $j$ & $(S,B),(C_j,C_j),(C_j,S),(C_{3-j},C_{3-j}),(B,B)$\\
\bottomrule
\end{tabular}
\end{center}
The row conditions are Boolean circuits on the finite codes. The
conditional requirements are covered by Lemma~\ref{lem:compiler}.
There are no existential row requirements: all semantic conditions
are preserved on subteams.

\section{Extracting a halting computation from a counterteam}
\label{sec:soundness}
Suppose a finite semantic model $R$ violates the core target $D$.
Pass to the subteam $S$ supplied by Lemma~\ref{lem:hall}; all semantic
conditions remain true. A degree-five token must use an $e_i$-edge,
since there are only four ordinary edge types and each type occurs
at most once at a token. The degree-five tokens are therefore exactly
the globally unique start and halt control tokens. Both have all five
allowed ports $a,c,e_1,e_2,e_3$. Every other token has exactly the
four ports $a,b,c,d$.

Follow the first channel forward from the start control token, taking
an $a$-edge at a $U$-token and a $b$-edge at a $V$-token. The kind
remains $S$. Every internal control token has exactly one incoming
and one outgoing edge in this directed channel. The start has no
incoming edge and the halt has no outgoing one. The walk cannot
enter a cycle: at the first repeated token, either the start acquires
an incoming edge or another token acquires two different incoming
edges of the same type. Finiteness therefore forces the walk to end
at halt. Reset edges cannot occur on it, since both their endpoints
have kind $B$.

The encountered configurations are distinct, since each has at most
one control token and the walk does not repeat a token. Each of these
configurations has exactly one control token. We claim that the
second channel between consecutive configurations is a bijection on
\emph{their entire token sets}. For a transition from a $U$-configuration
$X_0$ to a $V$-configuration $Y_0$, the control $a$-edge fixes the
neighbour pair at direction $0$. Every token of $X_0$ has a $c$-port,
and $\dep{X\delta}Y$ sends it to $Y_0$. Conversely, every token of
$Y_0$ has a $c$-port, and $\dep{Y\delta}X$ makes its predecessor
belong to $X_0$. Port uniqueness gives a bijection. The argument for
a transition in the other parity uses the $d$-ports and direction $1$.
The start and halt controls possess the required $c$-ports; every
other token possesses all ordinary ports.

There is one $S$ on either side of each bijection. For an increment
of counter $j$, the source $S$ becomes $C_j$ and the unique target
$S$ comes from one blank. Thus the number of $C_j$-tokens increases
by exactly one, while the other counter is unchanged. For a decrement,
the source $S$ becomes blank, and the target $S$ requires one source
$C_j$; hence the counter was positive and decreases by exactly one.
The no-op preserves both counters. On a zero-test transition, all
source tokens participate in the second channel, and its row guard
excludes every $C_j$-token. The test is therefore genuinely zero.
The control walk is a valid machine computation from zero counters
to halt.

\section{Constructing a counterteam from a halting computation}
\label{sec:completeness}
Conversely, take a halting computation of the normalized machine.
Choose a positive integer $N$ strictly larger than the maximum sum
of the counter values along it. Represent each configuration by $N$
tokens: one $S$, the required numbers of $C_1,C_2$, and the remaining
blanks. Give every configuration a fresh $X$- or $Y$-identifier
according to parity, and every token a fresh $U$- or $V$-identifier.

For each ordinary transition choose two independent bijections.
The first maps $S$ to $S$ and pairs the remaining tokens arbitrarily.
The second follows the operation table. Such a bijection exists:
for an increment use the source control token for the additional
counter unit and one blank for the target control token; for a
decrement use one counter unit for the target control token and
the source control token for the additional blank. Match all remaining
kinds identically. A zero test or no-op uses the identity on kinds.
The choice of $N$ supplies the blank needed by every increment.
Assign edge types $a,c$ or $b,d$ according to the transition parity.

Both endpoint configurations have zero counters and hence $N-1$
blanks. Join their blanks by two reset bijections of types $b,d$
from halt to start. Add three edges of types $e_1,e_2,e_3$ between
the start and halt control tokens. Give every edge a different
$w$-value.

The atoms $\dep{X\delta}Y$ and $\dep{Y\delta}X$ hold because each configuration has
one neighbour per direction, including reset and $e_i$ directions.
All other consistency and row conditions hold by construction.
Every token has degree four except the two degree-five endpoint
controls. To satisfy each degree premise, assign one $e_i$-edge
to its constancy part; after removing that edge, all degrees on
both sides are at most four. If the total number of tokens is $n$,
there are $2n+1$ edges. They cannot be assigned to the $2n$ available
endpoint slots, so $D$ fails.

\section{The reduction in the required syntax}
\label{sec:reduction}
Apply Lemma~\ref{lem:compiler} to the semantic theory of
Section~\ref{sec:machine}. Let $A$ contain all semantic variables,
including the Boolean guards, and call the resulting pure theory
$\Sigma_M$. Its only potentially unbounded variables are
$U,V,X,Y,w$. The marker $z$, all code bits, guards, switches, and
graph variables have explicit bounds of at most three values.
The key premises are $\dep w a$ for $a\in A$.

Given a halting computation, Section~\ref{sec:completeness} constructs
a finite semantic model violating $D$. Its canonical expansion
satisfies $\Sigma_M$, and Lemma~\ref{lem:filter} makes it violate
$(2\cdot\dep Az)\vee D$.
Conversely, a finite model of $\Sigma_M$ violating this filtered
target has a nonempty core violating $D$. By compilation and
injective Boolean decoding the core satisfies all semantic
conditions. Section~\ref{sec:soundness} extracts a halting
computation from it.

Since $w\in A$ and $\dep w a\in\Sigma_M$ for all $a\in A$, agreement on $w$
is exactly agreement on all of $A$, also on every subteam. Therefore
$\dep Az$ and $\dep wz$ are interchangeable throughout a cover.
The filtered target is equivalent, on all models of $\Sigma_M$,
to the fixed formula $\varphi_*$. We have proved
\[
 M\text{ halts}
 \quad\Longleftrightarrow\quad
 \Sigma_M\not\models_{\fin}\varphi_*.
\]
This is an effective many-one reduction, so nonimplication is
$\RE$-hard. Lemma~\ref{lem:enumeration} gives $\RE$-membership.
Complementation proves $\coRE$-completeness, and
Corollary~\ref{cor:finite} transfers the result to unrestricted teams.

The only five-occurrence premises are $P_U$ and $P_V$. Conditional
dependence requirements use two occurrences, and the remaining compiler premises use
at most three. All the explicit three-value bounds are themselves
three-occurrence disjunctions. This proves every restriction in
Theorem~\ref{thm:main}.

\begin{corollary}[No computable counterteam bound]
\label{cor:no-bound}
There is no computable function $B(\Sigma,\varphi)$ such that every
nonimplication instance has a counterteam of at most
$B(\Sigma,\varphi)$ rows. This remains true for the fixed consequent
and restricted premises of Theorem~\ref{thm:main}.
\end{corollary}
\begin{proof}
A computable bound would decide implication by enumerating the
finitely many equality types of teams up to that size and checking
all premises and the consequent. This contradicts
Theorem~\ref{thm:main}.
\end{proof}

Thus finite counterteams always suffice, but their sizes have no
computable uniform bound in terms of the instance. The reason is
not merely the existence of large minimal obstructions. Although
the premises alone are downward closed, failure of the consequent
supplies a Hall surplus. Combined with upper degree bounds, this
surplus forces every required port to exist and hence enforces an
entire computation of unbounded duration.

\begingroup
\small
\begin{thebibliography}{9}
\setlength{\itemsep}{3pt}
\bibitem{Vaananen07}
Jouko V\"a\"an\"anen.
\emph{Dependence Logic: A New Approach to Independence Friendly Logic}.
London Mathematical Society Student Texts 70,
Cambridge University Press, 2007.
\href{https://doi.org/10.1017/CBO9780511611193}{doi:10.1017/CBO9780511611193}.

\bibitem{Kontinen10}
Jarmo A. Kontinen.
\emph{Coherence and Complexity in Fragments of Dependence Logic}.
PhD thesis, Universiteit van Amsterdam, ILLC DS-2010-05, 2010.
\href{https://pure.uva.nl/ws/files/817636/77418_thesis.pdf}{Thesis PDF}.

\bibitem{FKM26}
Nicolas Fr\"ohlich, Phokion G. Kolaitis, and Arne Meier.
\emph{Disjunctions of Two Dependence Atoms}.
In CSL 2026, LIPIcs 363, 10:1--10:21, 2026.
\href{https://doi.org/10.4230/LIPIcs.CSL.2026.10}{doi:10.4230/LIPIcs.CSL.2026.10}.

\bibitem{Dudenhefner22}
{\raggedright
Andrej Dudenhefner.
\emph{Constructive Many-One Reduction from the Halting Problem to
Semi-unification}. In CSL 2022, LIPIcs 216, 18:1--18:19, 2022.
\href{https://doi.org/10.4230/LIPIcs.CSL.2022.18}{doi:10.4230/LIPIcs.CSL.2022.18}.\par}

\bibitem{Hall35}
Philip Hall. \emph{On Representatives of Subsets}.
Journal of the London Mathematical Society s1-10(1), 26--30, 1935.
\href{https://doi.org/10.1112/jlms/s1-10.37.26}{doi:10.1112/jlms/s1-10.37.26}.
\end{thebibliography}
\endgroup
\end{document}
